\section{Introduction}

Within-word variation in human speech poses challenges for ASR systems~\cite{radford2022whisper, zhang2023google-usm, pratap2024scaling-mms, puvvada2024less-canary-nvidia, nvidia2025parakeet} that focus solely on word-level transcription. Yet, modeling subword structures like syllables or phonemes is essential for applications in language learning and clinical analysis~\cite{kheir-etal-2023-automatic-assessment-review,truong04_icall-acoustuc-phonetic,lee2016language-capt,cordella2024connected-fluency1,fontan2023automatically-fluency2,gordon1998fluency-fluency3,gordon2022clinicians-fluency4,metu2023evaluating-fluency5,ssdm, lian2024ssdm2.0, zhou2024yolostutterendtoendregionwisespeech, zhou2024stutter, zhou2024timetokensbenchmarkingendtoend}. This work focuses on phoneme transcription—a particularly challenging yet critical task.

Several factors contribute to the difficulty of phoneme-level transcription. First, unlike word-level targets, phoneme annotations are inherently non-deterministic due to accent, allophony and other context-dependent realizations~\cite{choi2025leveragingallophonyselfsupervisedspeech}. Second, existing training objectives, such as CTC~\cite{graves2006connectionist-ctc} and attention-based models~\cite{chorowski2015attention-asr}, often introduce alignment noise. For example, CTC tends to produce overly peaked label distributions~\cite{zeyer2021does-ctc-peaky}, which result in inaccurate timing and error-prone alignments. Attention-based models, on the other hand, may hallucinate transcriptions~\cite{koenecke2024careless}—an issue that is particularly problematic in clinical applications where precision is critical. Third, most existing remedies are designed for fluent speech and demonstrate limited effectiveness on non-fluent speech. For example, several approaches~\cite{tian2022bayes-ctc1, huang2024less-ctc2, yao2024cr-ctc3} enhance CTC-based alignment by introducing heuristic constraints on the prior. However, these heuristic alignments have been shown to be incompatible with non-fluent speech~\cite{lian2023unconstrained-udm, lian-anumanchipalli-2024-towards-hudm,ssdm, lian2024ssdm2.0, ye2025seamlessalignment-neurallcs}. 

Given the aforementioned challenges, a key question arises: \textit{Can we design a phoneme transcription objective that performs robustly simultaneously on both fluent and non-fluent speech, while also producing reliable and accurately aligned outputs suitable for clinical applications?} 

A key insight arises from treating \textit{phoneme transcription as a speech production process}, particularly when speakers are prompted to read reference text. Fluent speech results from accurate production, while deviations introduce non-fluent speech. For example, when given the phoneme sequence “IH N S ER T,” a speaker might produce “IH [S] N S ER [AH] T,” introducing insertions on “IH” and “ER.” This yields an alignment such as “IH–(IH,S), N–(N), S–(S), ER–(ER,AH), T–(T)” (label–spoken). The core idea is that the \textit{loss function should account only for effective speech-text alignments}. In the case of “IH–(IH,S),” frames aligned to the extraneous “S” should be excluded when computing the loss for “IH.” Without such errors, the method reduces to the standard ASR objective. This alignment strategy follows the Longest Common Subsequence (LCS) principle~\cite{hirschberg1977algorithms-lcs1}, which was recently shown effective for non-fluent speech~\cite{ssdm}.
Building on this intuition, we apply the LCS algorithm online to identify effective alignments, which are then used to regularize the vanilla CTC~\cite{graves2006connectionist-ctc} paths, as shown in Fig.~\ref{fig:pipeline}. We refer to this method as \textbf{LCS-CTC}. Unlike previous methods that introduce heuristics to regularize CTC alignments~\cite{tian2022bayes-ctc1, huang2024less-ctc2, yao2024cr-ctc3}, \textit{LCS-CTC is grounded in human speech production, interpretable, and clinically meaningful}.

Our LCS-CTC framework consists of two main components: a phoneme-aware alignment module based on a modified Longest Common Subsequence (LCS) algorithm, and a constrained CTC training objective that incorporates the resulting alignment masks. The alignment module first predicts a frame-phoneme cost matrix using a lightweight neural network trained with weak supervision. This matrix encodes pairwise acoustic-phonetic similarity and is used to derive a partial alignment path via a similarity-aware LCS dynamic programming procedure. The resulting alignment mask identifies high-confidence frame-phoneme pairs and is used in two complementary ways during training: (1) it anchors specific frames to phoneme labels with cross-entropy supervision, and (2) it restricts the CTC path space via masked emission probabilities. Together, these two components enable robust phoneme recognition by combining interpretable alignment guidance with flexible sequence modeling.


We evaluate our method on both fluent (LibriSpeech~\cite{7178964}) and non-fluent corpora, including PPA speech~\cite{gorno2011classification-ppa} and the largest existing simulated non-fluent corpus, \textit{LLM\_dys}~\cite{zhang2025analysisevaluationsyntheticdata}. Results show that LCS-CTC consistently outperforms vanilla CTC across all metrics, including (weighted) phoneme error rate and duration-aware alignment accuracy. Notably, while LCS-CTC is inspired by non-fluent speech, it also provides stronger regularization for fluent speech. Although we have not yet explored its scalability across data types, linguistic diversity, or domains, the method shows strong potential and may serve as a foundation for universal phoneme recognition. We have open-sourced our model and checkpoints at \href{https://github.com/Auroraaa86/LCS-CTC}{\texttt{https://github.com/Auroraaa86/LCS-CTC}}










